\documentclass[11pt]{article}
\usepackage[a4paper,margin=25mm]{geometry}
\usepackage{amsmath,amssymb,amsthm,mathrsfs,hyperref}
\setlength{\emergencystretch}{3em}
\newtheorem{theorem}{Theorem}
\newtheorem{lemma}[theorem]{Lemma}
\newtheorem{proposition}[theorem]{Proposition}
\newcommand{\MA}{\operatorname{MA}}
\newcommand{\Ent}{\operatorname{Ent}}
\newcommand{\Ric}{\operatorname{Ric}}
\newcommand{\NA}{\mathrm{NA}}
\newcommand{\PSH}{\operatorname{PSH}}
\title{Joint entropy recovery at a prescribed finite-entropy start\\
\large An independently checked supporting theorem for P19.2}
\author{Scoped theorem; no claim of a complete solution of P19.2}
\date{6 October 2026}
\begin{document}
\maketitle
\begin{abstract}
Working with the four complete mathematical grants in the input, we
prove a joint entropy/time recovery theorem for maximal weak rays on
every fixed smooth compact K\"ahler manifold. The actual finite-entropy
initial potential is retained and can be unbounded and nonmodel.
Strong convergence of non-Archimedean data, entropy recovery, finite-time
lower semicontinuity, and recession slopes suffice: no uniform mixed
pairing error estimate is needed. A convex chord argument treats the
cohomological correction as a difference of convex contracted energies.
There is also a joint divergence theorem for infinite non-Archimedean
entropy. These conclusions improve the general fixed-start mechanism
recorded in grant S, but construct neither closed relative P-profile
recovery nor approximations of arbitrary nonmaximal rays. The entire
original P19.2 is not solved, and no full-scope analytic obstruction
is proved. The results have been independently checked relative to
the stated complete caller grants.
\end{abstract}

\section{The unchanged task and the result obtained}
The exact original P19.2 is: construct model approximations in the
residual hybrid programme preserving all required energy, entropy and
singularity data; energy convergence alone is insufficient, and the
published/source-stated projective Fujita regimes are established
subcases. The full supplied assignment additionally retains arbitrary
actual finite-energy weak rays, their actual prescribed singular starts,
closed finer P data, and raw mixed masses, including genuine zero-mass
and identically minus-infinite levels as distinct data. It requires
the original carrier with supported semivaluations, discrepancy factors,
actual positive models or flag configurations, general compact
K\"ahler/transcendental classes, and the joint entropy/time normalization.
It expressly does not permit an argument about one maximalization
mechanism to be counted as universal model-approximation nonexistence.

We preserve this task rather than replacing it with our theorem's
signature. The complete original input, task and assignment are frozen
as local artifacts, with the native notes list empty. Our theorem is
restricted to the maximal branch with a finite-entropy actual start.
It does not assert convergence of entropy at a fixed positive time,
or of any prescribed closed P-profile. The new claim is joint
entropy/time recovery in this general fixed-manifold maximal branch.
The original atom therefore remains unresolved.

\section{Conventions, carrier, and actual models}
Let $X$ be connected, smooth and compact K\"ahler, of dimension $n\ge1$.
Fix a K\"ahler form $\omega$, $\alpha=[\omega]$,
$V=\int_X\omega^n$ and $\mu_0=\omega^n/V$. Throughout,
\[
 dd^c=\frac{i}{2\pi}\partial\bar\partial,\qquad
 t=-\log|\tau|^2,\qquad dd^c\log|f|^2=[\operatorname{div}(f)].
 \tag{1}
\]
For $u\in\mathcal E^1(X,\omega)$ use the complete finite-energy
extensions of
\[
 E(u)=\frac1{(n+1)V}\sum_{k=0}^n
  \int_X u\,\langle\omega_u^k\rangle\wedge\omega^{n-k},\quad
 \MA(u)=V^{-1}\langle\omega_u^n\rangle,\quad
 H(u)=\Ent(\MA(u)\mid\mu_0).
 \tag{2}
\]
These are full probability conventions: $E(0)=0$, $E(u+c)=E(u)+c$,
and entropy is infinite without absolute continuity. Here
$\omega_u=\omega+dd^cu$, and brackets denote nonpluripolar products.
An actual weak ray is strongly continuous at time zero and is the
canonical finite-energy weak geodesic on every finite interval.
Its total potential is $U(x,\tau)=u_{-\log|\tau|^2}(x)$.

The non-Archimedean carrier is the original full $X^\beth$ in the
trivially valued theory of the grants. Values $+\infty$ on the support
of a semivaluation are retained. The usual datum is obtained from
the generic divisorial coefficients of the actual total potential,
after an affine time normalization where required. It is not assigned
by a pointwise quotient of an infinite value by time.

An actual relatively K\"ahler model in this paper is a smooth
equivariant proper modification
$\rho:\mathcal Y\to X\times\mathbf P^1$, trivial away from the central
fiber, with a vertical real divisor $D$ and relatively K\"ahler class
$\mathcal A=\rho_X^*\alpha+[D]$. The construction imported from R
produces such geometric models, not just formal PL functions.
In a general transcendental class this is cohomological data;
there need not be a rational line bundle. We do not call a merely
nef or log-discrepancy class semiample.

On a central Gauss component $F$, put $b_F=\operatorname{ord}_F(\tau)$.
The model datum is $f_D(v_F)=D_F/b_F$, with compatibility sign
$U+\log|s_D|^2=O(1)$ for bounded compatible model weights.
The component formulas are
\[
 \MA_\alpha(f_D)=\frac1V\sum_F b_F(\mathcal A^n\cdot F)\delta_{v_F},
 \qquad E_\alpha(f_D)=\frac{\mathcal A^{n+1}}{(n+1)V},
 \qquad H_\alpha(\psi)=\int_{X^\beth}A_X\MA_\alpha(\psi).
 \tag{3}
\]
The compactification in (3) is trivialized near infinity as in R.
The top measure has total mass one. On an SNC model the discrepancy
factor on $F$ is $(\kappa_F+1-b_F)/b_F$, where $\kappa_F$ is the
relative canonical coefficient. Thus neither nonreduced central
multiplicities nor discrepancy factors are discarded. Strong limits,
energies, and mixed pairings in (3) are those of grant R.

For a prescribed start $w$ write $r(w,\psi)$ for its actual maximal
ray. The model datum does not make this ray's residual weight bounded
when $w$ is unbounded. We keep $w$ exactly, rather than replacing it
with a model initial potential.

\section{Exact premises and new source application}
The complete caller bodies R, T, C and S, and only those bodies, are
grants. Their historical verdicts are not premises. The following
specific consequences of R are used.
\begin{enumerate}
\item For every actual $w\in\mathcal E^1$ and
$\psi\in\mathcal E^{1,\NA}(\alpha)$, $r(w,\psi)$ exists, is ordered
in both arguments, and has
$E(r(w,\psi)_t)=E(w)+tE_\alpha(\psi)$. Its usual datum is $\psi$.
Strong convergence $\psi_j\to\psi$ with the same $w$ gives bounded-time
$d_1$ convergence of the maximal rays.
\item For maximal rays with a finite-entropy slice,
$\lim_{t\to\infty}H(r_t)/t=H_\alpha(\psi)$, including $+\infty$.
When $H(w)$ and $H_\alpha(\psi)$ are finite, every finite-time
entropy is finite. An arbitrary nonmaximal actual ray has infinite
entropy slope, and the original nonmaximal ray has positive energy defect relative to its
maximalization.
\item Fixed smooth cohomological contracted energies have calibrated,
finite recession slopes on maximal rays. These slopes are strongly
continuous in $\psi$. Energy and the standard finite-energy
comparison/decreasing-geodesic foundations are also granted.
\item When $H_\alpha(\psi)<\infty$, there are actual relatively
K\"ahler model data $\eta_j\to\psi$ strongly with
$H_\alpha(\eta_j)\to H_\alpha(\psi)$ (R, equation (16)).
This construction does not grant finer P-profile recovery.
\end{enumerate}
We do not independently reprove R or its external theorem chain.
T supplies its zero-start one-branch $\mathbf P^1$ result only;
C its finite-entropy absolute-endpoint example on $\mathbf P^n$,
$n\ge2$, only; S its scoped singular-start models, relative lemma,
closed-level calculations and stated general mechanism only. These
three are used below to delimit scope, not to expand the main theorem.
The ledger supplied inside S is historical source data and has not
been treated as our own reading record.

One primary source is newly applied: the published
Berman--Darvas--Lu paper \cite{BDL}. On a connected compact K\"ahler
manifold its Proposition 4.4 says that the contracted energy
$\mathrm{AM}_\gamma$, for a smooth closed real $(1,1)$ form $\gamma$,
extends $d_1$-continuously to $\mathcal E^1$ and is bounded on
$d_1$-bounded sets. Proposition 4.5 gives its convexity for smooth
positive $\gamma$ along weak geodesics with bounded Laplacian.
Theorem 4.7, at twist zero and $p=1$, gives $d_1$ lower
semicontinuity and geodesic convexity of the extended Mabuchi functional.
The complete statement/proof units for these applications were read;
literal source coordinates and precise convention conversion are
recorded in the appendix. No projectivity occurs in these statements.

The source uses $i\partial\bar\partial$. Set $v=u/(2\pi)$ with the
same background $\omega$; the currents and probability measures coincide.
Its distance is $d_{1,\mathrm{BDL}}=d_1/(2\pi)$,
$\mathrm{AM}(v)=E(u)/(2\pi)$, and for the same smooth form $\gamma$,
\[
 T_\gamma(u):=\frac1V\sum_{k=0}^{n-1}
  \int_Xu\,\gamma\wedge\langle\omega_u^k\rangle
                     \wedge\omega^{n-1-k}
 =2\pi n\,\mathrm{AM}_\gamma(v).                    \tag{4}
\]
Consequently these continuity and convexity assertions keep their
sign and domain. To extend the bounded-geodesic convexity in
Proposition 4.5 to an arbitrary finite-energy segment, decrease smooth
endpoints to its endpoints. The canonical geodesics decrease and
converge in $d_1$ at each time by the granted R foundations.
Apply Proposition 4.4 at each time to their contracted energies.
Taking the limit of the convexity inequalities proves convexity
of $T_\gamma$ for the desired segment. This argument explicitly
uses the finite-energy extension and does not extrapolate the
bounded statement without proof.

Put $\Ric(\omega)=-dd^c\log\omega^n$ and
$\overline S=nV^{-1}\int_X\Ric(\omega)\wedge\omega^{n-1}$. Define
\[
 M(u)=H(u)+G(u),\qquad
 G(u)=T_{-\Ric(\omega)}(u)+\overline S E(u).           \tag{5}
\]
The source's Ricci form and mean scalar are $2\pi$ times ours.
Using (4) in its untwisted formula gives exactly (5), with entropy
coefficient one. Thus $M$ is lower semicontinuous and convex in
our convention. This fixes both probability and squared-log time
normalizations rather than importing an uncalibrated slope.

\section{The convex chord argument}
\begin{lemma}[Finite recession recovery]\label{convex}
Let $f_j:[0,\infty)\to\mathbf R$ be finite convex functions with
$f_j(0)\to a\in\mathbf R$. Let $f$ be finite convex with $f(0)=a$.
Suppose for every $T>0$, $\liminf_j f_j(T)\ge f(T)$,
and
\[
 \ell_j=\lim_{t\to\infty}f_j(t)/t\to\ell
           =\lim_{t\to\infty}f(t)/t\in\mathbf R.
\]
Then $f_j(t_j)/t_j\to\ell$ for every $j\to\infty$ and
$t_j\to\infty$.
\end{lemma}
\begin{proof}
Convexity makes $(f_j(t)-f_j(0))/t$ nondecreasing and bounded above
by $\ell_j$. For $t\ge T>0$,
\[
 \frac{f_j(0)}t+\frac{f_j(T)-f_j(0)}T
 \le\frac{f_j(t)}t\le\frac{f_j(0)}t+\ell_j.          \tag{6}
\]
Along any joint sequence the upper bound tends to $\ell$ and the
lower bound has liminf at least $(f(T)-a)/T$. Let $T\to\infty$.
The two bounds coincide. The same proof applies to any subsequence
of indices, so the asserted quantifier is genuinely joint.
\end{proof}

\begin{lemma}[Infinite recession lower recovery]\label{infinite}
Keep $f_j$ finite convex with $f_j(0)\to a$.
Let $f:[0,\infty)\to\mathbf R\cup\{+\infty\}$ be extended convex,
with $f(0)=a$, $\liminf_j f_j(T)\ge f(T)$ at every $T>0$, and
$\lim_{t\to\infty}f(t)/t=+\infty$.
Then $f_j(t_j)/t_j\to+\infty$ for every joint sequence.
No convergence assumption on $\ell_j$ is needed.
\end{lemma}
\begin{proof}
Use the lower bound in (6). If $f(T)=+\infty$ at some fixed $T$,
it already implies divergence along the joint sequence. Otherwise
all $f(T)$ are finite, and $(f(T)-a)/T$ is arbitrarily large for
large $T$. First fix such a $T$, take the joint liminf, and then
let $T$ increase. This gives divergence to $+\infty$.
\end{proof}

Neither lemma proves $f_j(T)\to f(T)$. The finite lemma uses only
a one-sided inequality at a fixed time. It also gives no uniform
additive $O(1)$ recession remainder.

\section{Joint cohomological slopes}
For a smooth closed form $\gamma$ write $L_\gamma(\psi)$ for the
calibrated recession slope of $T_\gamma(r(w,\psi)_t)$ supplied by R.
It is finite, independent of the chosen finite-energy start, and
strongly continuous in $\psi$. For a model $f_D$ its exact
intersection expression is
\[
 L_\gamma(f_D)=\frac{\mathcal A^n\cdot\rho_X^*[\gamma]}V.
 \tag{7}
\]
This is the contracted mixed-energy slope, not a probability
renormalization of a raw level mass.

\begin{proposition}\label{cohomological}
Let $\psi_j\to\psi$ strongly in $\mathcal E^{1,\NA}(\alpha)$,
keep any $w\in\mathcal E^1$ fixed, and put
$r_j=r(w,\psi_j)$, $r=r(w,\psi)$.
For every smooth closed real $(1,1)$ form $\gamma$ and every joint
$j\to\infty$, $t_j\to\infty$,
\[
 \frac{T_\gamma(r_{j,t_j})}{t_j}\longrightarrow L_\gamma(\psi),
 \qquad
 \frac{G(r_{j,t_j})}{t_j}\longrightarrow G_\alpha(\psi),\quad
 G_\alpha=L_{-\Ric(\omega)}+\overline S E_\alpha.      \tag{8}
\]
\end{proposition}
\begin{proof}
Assume first that $\gamma\ge0$. Its contracted energy is a finite
convex function along each finite-energy ray, by Section 3.
At any fixed $T$, R gives $r_{j,T}\to r_T$ in $d_1$.
Proposition 4.4 after (4) gives
$T_\gamma(r_{j,T})\to T_\gamma(r_T)$.
Initial values are all $T_\gamma(w)$. Their finite recession
slopes converge by R's strongly continuous calibrated pairing.
Lemma \ref{convex} therefore proves the first limit.

For any signed $\gamma$ choose a fixed $c>0$ sufficiently large
that $\gamma+c\omega\ge0$. Write
$\gamma=(\gamma+c\omega)-c\omega$ and apply the preceding
argument to both terms. Linearity proves the assertion for
signed forms. This decomposition is fixed in $j$ and $t$.
Finally, energy is affine:
$E(r_{j,t})/t=E(w)/t+E_\alpha(\psi_j)$.
Strong energy convergence and (5) give the second limit.
No estimate depending on a time-varying $d_1$ ball or an unnormalized
relative $J$ is used.
\end{proof}

\section{The new maximal-branch entropy theorem}
\begin{theorem}[Joint entropy recovery with the actual start retained]
\label{entropy}
Fix the general smooth compact K\"ahler manifold and conventions
of Section 2. Let $w\in\mathcal E^1$ satisfy $H(w)<\infty$.
Suppose $\psi_j\to\psi$ strongly in
$\mathcal E^{1,\NA}(\alpha)$, with $H_\alpha(\psi_j)<\infty$.
Set $r_j=r(w,\psi_j)$ and $r=r(w,\psi)$.
\begin{enumerate}
\item If $H_\alpha(\psi)<\infty$ and
$H_\alpha(\psi_j)\to H_\alpha(\psi)$, then for every joint
$j\to\infty$, $t_j\to\infty$,
\[
 \frac{H(r_{j,t_j})}{t_j}\longrightarrow H_\alpha(\psi).
 \tag{9}
\]
Equivalently, for every $\epsilon>0$ there exist $J,T$ such that
$|H(r_{j,t})/t-H_\alpha(\psi)|<\epsilon$ for all $j\ge J,t\ge T$.
\item If $H_\alpha(\psi)=+\infty$, then every joint sequence
satisfies $H(r_{j,t_j})/t_j\to+\infty$, without an additional
assumption on $H_\alpha(\psi_j)$'s limit.
\end{enumerate}
In both cases energy recession is recovered jointly, and
$r_{j,t}\to r_t$ in $d_1$ uniformly on bounded time intervals.
The start is exactly $w$, including when it is unbounded and nonmodel.
\end{theorem}
\begin{proof}
For each $j$, grant R gives finite entropy at every finite time,
the entropy recession $H_\alpha(\psi_j)$, and calibrated
cohomological recession $G_\alpha(\psi_j)$.
Thus $f_j(t):=M(r_{j,t})$ is finite convex and has finite slope
\[
 m_j=H_\alpha(\psi_j)+G_\alpha(\psi_j),\qquad f_j(0)=M(w)<\infty.
 \tag{10}
\]
At fixed $T$, bounded-time recovery and lower semicontinuity give
\[
 M(r_T)\le\liminf_j M(r_{j,T}).                         \tag{11}
\]
In the finite branch, R gives finite-time entropy for $r$ and
$\lim_{t\to\infty}M(r_t)/t
=m:=H_\alpha(\psi)+G_\alpha(\psi)$.
Entropy recovery and strong mixed-slope continuity give $m_j\to m$.
Apply Lemma \ref{convex} to (10)--(11). This proves the joint limit
of $M(r_{j,t_j})/t_j$. Subtract the independent joint limit (8).
The result is (9), with entropy coefficient one.

In the infinite branch, R's entropy classification and finite
cohomological recession give $M(r_t)/t\to+\infty$ as an extended
limit. $M(r_t)$ might be infinite at some finite times, which is
allowed in Lemma \ref{infinite}. Equation (11) and that lemma
give $M(r_{j,t_j})/t_j\to+\infty$. The joint correction (8) is
finite, so the same conclusion holds for entropy.

The energy assertion is the exact identity
\[
 \frac{E(r_{j,t_j})-E(w)}{t_j}=E_\alpha(\psi_j)\to E_\alpha(\psi).
 \tag{12}
\]
Bounded-time $d_1$ recovery is the granted R maximal-ray continuity.
The uniform-tail formulation in the finite branch is equivalent
to (9): failure would select $j_k\ge k,t_k\ge k$ violating it.
This completes the proof.
\end{proof}

\begin{proposition}[Actual model-sequence consequence]\label{models}
For every fixed actual $w\in\mathcal E^1$ with $H(w)<\infty$ and
$\psi\in\mathcal E^{1,\NA}(\alpha)$ with finite entropy,
there are actual relatively K\"ahler model data $\eta_j$ for which
the maximal rays from this unchanged $w$ have bounded-time $d_1$
recovery, energy-slope recovery and every joint entropy/time limit
in Theorem \ref{entropy}. No projectivity is assumed.
\end{proposition}
\begin{proof}
Take the exact model recovery in grant R, equation (16), retaining
its real cohomological model meaning. It gives strong convergence
and NA entropy convergence. Apply Theorem \ref{entropy}.
This application adds a joint analytic conclusion to that recovery;
it does not add a closed relative boundary theorem to R's model
construction. The model starts are not asserted to be model functions.
\end{proof}

The theorem closes the general joint-entropy gap of the specific
finite-entropy maximal recovery mechanism recorded in S, Section 5.
It does so by convexity rather than the mixed-pairing normalization
argument that S left unproved. No $O(1)$ uniform remainder is asserted.
This conclusion is separate from S's explicit radial formula, which
has the stronger uniform additive remainder in its precise scope.

\section{What the proof does not recover}
For a ray from its actual $w$, define the closed Legendre levels and
relative ceiling by
\[
 \chi_u(\lambda)=\inf_{t\ge0}(u_t-t\lambda),\qquad
 \mathscr P_w[h]=\left(\lim_{C\to\infty}P_\omega(h+C,w)\right)^*.
 \tag{13}
\]
Grant S proves recognition $\mathscr P_w[\chi_u(\lambda)]
=\chi_u(\lambda)$ whenever this is a genuine psh level. It uses
Perron comparison, not NA maximality. The case $-\infty$ is only
a placeholder convention. A ceiling against $w$ is distinct from
an absolute ceiling against zero. An I-envelope uses multiplier-ideal
data and cannot be substituted for this finer P-envelope.

For genuine levels use the raw mixed masses
\[
 M_k(h)=\frac1V\int_X\omega^{n-k}\wedge
                 \langle\omega_h^k\rangle,\qquad 1\le k\le n.
 \tag{14}
\]
No probability division by a zero mass is performed. A genuine
zero-mass level is still a potential and can have finer P data;
an identically minus-infinite placeholder has no current or
singularity metric. Only the explicit zero extension of the
positive-degree numbers (14), where stated, is used at that boundary.

The strong model sequence of Proposition \ref{models} need not be
ordered in a way controlling (13), and the argument uses only one
fixed time at a time. It supplies no interchange of an index limit
with the infimum over the entire half-line. The recognition lemma
is not a density or convergence theorem for relative ceilings.
Even when decreasing rays provide pointwise decreasing levels,
that statement alone is not a mixed nonpluripolar mass continuity
theorem at arbitrary singularities.

The grant boundaries illustrate these gaps accurately. T controls
all closed levels in its finite-height, zero-start, one-branch curve
domain. Its bad roof sequence for the trivial target diagnoses a
failed implication; the product sequence still recovers that target.
C changes an absolute zero-mass endpoint to $-\infty$ under a
finite-entropy singular starting-point change on $\mathbf P^n$,
$n\ge2$. It does not obstruct relative recovery retaining $w$.
S constructs actual rational relatively ample hyperplane blowup
models retaining one such unbounded start. Its genuine relative
endpoints descend to $-\infty$, with all positive-degree raw
mixed masses tending to zero, and its entropy estimate is joint.
That finite-dimensional radial family does not approximate arbitrary
data or arbitrary starts. The present theorem extends its general
joint-entropy mechanism, not its closed relative P-profile formula.

Finally, let $u$ be an actual nonmaximal ray from $w$ with usual
datum $\psi$, energy slope $c_u$, and
$\delta=E_\alpha(\psi)-c_u>0$. Grant R gives
\[
 d_1(r(w,\psi)_t,u_t)=t\delta.                         \tag{15}
\]
Strongly converging maximal model data in Proposition \ref{models}
converge at bounded times to $r(w,\psi)$, not $u$; energy's
$1$-Lipschitz property also gives, at each $j$,
\[
 \liminf_{t\to\infty}\frac{d_1(r(w,\eta_j)_t,u_t)}t
 \ge |E_\alpha(\eta_j)-c_u|\longrightarrow\delta.
 \tag{16}
\]
This is exactly a limitation of this maximal-ray mechanism.
It is not a proof that every permissible model or flag approximation
fails. We neither delete the arbitrary-ray branch nor manufacture a
universal obstruction from (16).

\section{Precise remaining gaps and comparison with the catalogue}
To settle the entire unchanged original P19.2, still missing here are:
\begin{enumerate}
\item A simultaneous model construction for arbitrary prescribed
finer P data and arbitrary finite-energy starts, with closed
relative levels retained on general compact K\"ahler manifolds.
\item General convergence of those actual relative ceilings and
all raw mixed masses in the required closed-level topology,
distinguishing genuine zero-mass data from $-\infty$ placeholders.
\item An approximation mechanism for arbitrary nonmaximal actual
weak rays preserving the actual ray and energy slope, rather than
replacing it by its maximalization.
\item The entropy requirements outside this finite-initial-entropy
maximal signature, including the requested fixed-time entropy
recovery. Strong $d_1$ convergence was used for lower
semicontinuity only, and proves no entropy continuity.
\item Residual log-discrepancy-model regularization in the required
generality and with the additional P data, with actual positivity
proved and no unjustified relative semiampleness assumption.
\end{enumerate}
The catalogue's published/source-stated smooth polarized special
Fujita regimes and the assignment's zero-start polarized
log-discrepancy-model recovery remain established subcases.
They do not supply the missing items by description alone.
Our reliance on the caller R grant allows a general smooth compact
K\"ahler maximal theorem without replacing the class by a polarization.
It does not resolve the original residual scope. P19.1 is granted,
and P19.3/P19.4 remain separate future targets; none is claimed here.
No rigorous full-scope analytic counterexample has been obtained.
Accordingly the requested completion flag is \texttt{candidate=false}.

\appendix
\section{Literal source coordinates and reproducible record}
The full supplied bodies were read in the input and extracted exactly
from its assignment; their literal coordinates in this run are
line numbers in the extracted files. The applications of R are
its Section 3 calibrated mixed slopes, equation (16), Sections 6--9
maximal-ray continuity, entropy classification and defect.
The relevant explicit assertions occur at local R lines 199--201,
467--507, 511--543 and 551--654. S's fixed-start recovery and
mechanism limitations are local S lines 459--507. These are fresh
local coordinates and are distinct from historical source layouts.

The primary BDL source is the published 44-page PDF,
Geometry \& Topology 21 (2017), 2945--2988, DOI
10.2140/gt.2017.21.2945, at the URL in the bibliography.
Returned layout lines were retained verbatim in JSON source records.
The complete units used here are Proposition 4.4 statement/proof
1852--1938 (context and product identification 1778--1948),
Proposition 4.5 statement/proof 2004--2085, and Theorem 4.7
statement/proof 2165--2307. Twist condition (3) was read at
490--506; untwisted and contracted conventions at 79--127,
1784--1834 and 508--532. The source ledger lists all returned
ranges, rather than claiming to have read the entire PDF.
The new applications use only smooth forms, twist zero and $p=1$;
their recursive pluripotential foundations are supplied by the
explicit caller R grant where invoked. There is no new application
to singular twisting currents or a degenerating family.

The unchanged native author records and reproducibility files are retained
in the round155 artifact directory; its public partial reader preserves
the original eight-page manuscript. Independent native Xean review156
read the whole manuscript and passed its precisely stated scope relative
to the four complete caller grants. This final paper corrects only the
harmless sentence that assigned the positive energy defect to the
maximalization rather than to the original nonmaximal ray. The equations
and all mathematical arguments are unchanged. The publication binding
and quote-free source application ledger record exact source and review
hashes and reading scopes. The original P19.2 remains open.

\begin{thebibliography}{9}
\bibitem{R} Caller grant R, complete supplied mathematical body,
\emph{Entropy slopes of finite-energy K\"ahler rays: calibration,
recovery, and arbitrary starting points}, 6 October 2026.
\bibitem{T} Caller grant T, complete supplied mathematical body,
\emph{Preserving the closed singularity profile in actual toric
model approximation}, 6 October 2026; exact zero-start curve scope.
\bibitem{C} Caller grant C, complete supplied mathematical body,
\emph{A finite-entropy initial singularity can erase an absolute
zero-mass endpoint}, 6 October 2026; $\mathbf P^n$, $n\ge2$.
\bibitem{S} Caller grant S, complete supplied mathematical body,
\emph{Retaining a singular start in model approximation:
relative P ceilings and an explicit collapsing boundary},
6 October 2026; explicit radial scope and general stated lemmas.
\bibitem{BDL} R. Berman, T. Darvas and C. H. Lu,
\emph{Convexity of the extended K-energy and the large time behavior
of the weak Calabi flow}, Geometry \& Topology 21 (2017), 2945--2988,
\href{https://msp.org/gt/2017/21-5/gt-v21-n5-p08-s.pdf}{published primary PDF},
Propositions 4.4--4.5 and Theorem 4.7.
\end{thebibliography}
\end{document}
